% Laget automatisk av elevark/lag_elevark.py fra opplegg/repetisjon-energi.html. Endringer her blir overskrevet.
\documentclass[11pt]{article}
\usepackage{elevark}
\begin{document}
\elevarkhode{Fysikk 1 · Mekanisk energi}{Repetisjon: arbeid og energi}{https://lektorbo.no/opplegg/repetisjon-energi.html}{Opplegget på nett}

Energi kan gå over fra en form til en annen, men forsvinner aldri. Med energibevaring kan vi ofte finne farten uten å vite noe om banen. Her er en oversikt og blandede oppgaver med løsningsforslag.

\section*{Når er den mekaniske energien bevart?}

\begin{center}\small
\begin{tabularx}{\linewidth}{YY}
\toprule
\textbf{Situasjon} & \textbf{Energiregnskap} \\
\midrule
Bare tyngden gjør arbeid (normalkraft og snorkraft vinkelrett på farten gjør ikke arbeid) & $E_{\mathrm{k}} + E_{\mathrm{p}}$ er konstant \\
Friksjon eller luftmotstand & $E_{k_{0}} + E_{p_{0}} = E_{\mathrm{k}} + E_{\mathrm{p}}$ + varme. Varmen er lik friksjonsarbeidet, Rs. \\
En ytre kraft gjør arbeid (motor, person) & Arbeidet fra kraften kommer i tillegg: $W = ΔE_{\mathrm{k}} + ΔE_{\mathrm{p}}$ + varme \\
\bottomrule
\end{tabularx}
\end{center}

\section*{Formler}

\begin{itemize}
\item Arbeid: $W = Fs \cos  α$, der $α$ er vinkelen mellom kraften og forflytningen.
\item Kinetisk energi: $E_{\mathrm{k}} = \tfrac{1}{2}mv^{2}$. Potensiell energi: $E_{\mathrm{p}} = mgh$.
\item Arbeid–energi: Det totale arbeidet på en gjenstand er lik endringen i kinetisk energi, $W_{\mathrm{tot}} = ΔE_{\mathrm{k}}$.
\item Effekt: $P = W/t = Fv$. Virkningsgrad: η = nyttig energi / tilført energi.
\end{itemize}

\section*{Oppgaver}

Bruk $g = 9{,}81\,\mathrm{m/s^{2}}$.

\begin{oppgaveliste}
\oppgave \fase{forklar}{Begrep}Er påstandene sanne eller usanne? Begrunn svaret.

a) Normalkraften gjør ikke arbeid på en vogn i en berg-og-dal-bane.\newline b) Dobbel fart gir dobbel kinetisk energi.\newline c) Uten friksjon får en kloss samme fart nederst i en bratt og en slak bakke med samme høyde.\newline d) Når du bærer en sekk bortover med konstant fart, gjør du positivt arbeid på sekken.\newline e) Energi som blir til varme, er tapt for alltid.
\svar
\oppgave \fase{test}{Regn}En ball kastes rett ned med 5,0~m/s fra 10~m høyde. Hvor fort treffer den bakken? Ville svaret vært annerledes hvis ballen ble kastet vannrett med 5,0~m/s?
\svar
\oppgave \fase{test}{Regn}En vogn i en berg-og-dal-bane har farten 2,0~m/s på toppen av en bakke 25~m over bakken. Hvor fort går den 8,0~m over bakken? I hvilken høyde har den farten 10~m/s? Se bort fra friksjon. Sjekk med simuleringen.\footnote{\url{lektorbo.no/sim/berg-og-dal.html}}
\svar
\oppgave \fase{test}{Regn}Et akebrett med barn har massen 40~kg og starter i ro på toppen av en 12~m høy bakke. Nederst har det farten 11~m/s. Hvor mye energi ble til varme? Bakken er 60~m lang. Hvor stor var den gjennomsnittlige friksjonskraften?
\svar
\oppgave \fase{test}{Regn}En heis på 600~kg løftes 20~m på 25~s med konstant fart. Finn arbeidet og effekten. Motoren har virkningsgrad 75~\%. Hvor stor effekt må motoren få tilført?
\svar
\oppgave \fase{test}{Regn}En sprettball slippes fra 1,80~m og spretter opp til 1,35~m. Hvor stor del av energien er igjen etter sprettet? Hvor høyt kommer ballen etter tre sprett hvis den mister like stor del hver gang? Sjekk med simuleringen.\footnote{\url{lektorbo.no/sim/sprettball.html}}
\svar
\oppgave \fase{utfordring}{Sammensatt}En elev på 65~kg løper opp en trapp med høyden 3,0~m på 4,0~s. Finn effekten mot tyngden. En elev hevder at hun har like stor effekt som en støvsuger på 500~W. Har hun rett? Hvorfor bruker kroppen likevel mye mer energi enn mgh? Se Trappeløp.\footnote{\url{lektorbo.no/sim/trappelop.html}}
\svar
\end{oppgaveliste}

\end{document}
